\SourceLeaf{531}
\begin{center}
\begin{tabular}{rrr@{\hspace{2em}}rrr}
\shortstack{Tausender-\\gruppe} & \textsc{Gauss} & \shortstack{Wahrer\\Wert} & \shortstack{Tausender-\\gruppe} & \textsc{Gauss} & \shortstack{Wahrer\\Wert} \\
$20$ & $102$ & $104$ & $546$ & $68$ & $69$ \\
$159$ & $87$ & $77$ & $601$ & $75$ & $76$ \\
$199$ & $96$ & $86$ & $625$ & $68$ & $78$ \\
$206$ & $85$ & $83$ & $668$ & $73$ & $74$ \\
$245$ & $78$ & $88$ & $675$ & $69$ & $73$ \\
$289$ & $85$ & $77$ & $784$ & $74$ & $75$ \\
$290$ & $84$ & $85$ & $800$ & $81$ & $71$ \\
$334$ & $80$ & $81$ & $879$ & $68$ & $78$ \\
$352$ & $80$ & $81$ & $985$ & $74$ & $70$ \\
$501$ & $78$ & $79$ &  &  &  \\
\end{tabular}
\end{center}
\par Die im Brief von \textsc{Gauss} an \textsc{Encke} erwähnte Formel \textsc{Encke}’s scheint die folgende
\[
\frac{n}{\log n}\overset{2\log n}{\surd}10
\]
zu sein, die \textsc{Encke} in einem Brief an \textsc{Gauss} vom 4. Dezember 1849 mitteilt.
\par Die \textit{Tafel der Anzahl der Klassen binärer quadratischer Formen} gibt die Anzahl der Geschlechter und Klassen sowie den Index der Irregularität für die negativen Determinanten in den Hundertergruppen 1 bis 30, 43, 51, 61, 62, 63, 91 bis 100, 117 bis 120 an; ferner in einer gesonderten Zusammenstellung für die der 1., 3. und 10. Tausendergruppe, für die ersten 800 der Formen $-(15n+7)$ und $-(15n+13)$ sowie für einige sehr große Determinanten; außerdem für die positiven Determinanten der 1., 2., 3., 9. und 10. Hundertergruppe und für einige weitere. Die Handschrift besteht aus einzelnen Zetteln mit unterschiedlich eingerichteten Tafeln. So steht bei den älteren das Wort Ordo anstelle von Genus, bei den einzelnen Hundertergruppen mit Ausnahme der 9. und 10. positiver Determinanten, aber auch bei einzelnen vorläufigen Zusammenstellungen in Tausendergruppen.
Zur besseren Übersicht ist hier durchweg die Bezeichnung der Disqu. Arithm. gewählt. Aus Platzmangel stehen die größten und kleinsten Quotienten der Klassenanzahl geteilt durch die Determinante sowie die Anzahl der Determinanten, deren Quotient innerhalb bestimmter Grenzen liegt, nicht unter den einzelnen Hundertergruppen, sondern zusammen am Ende der Tafel für die negativen Determinanten. Einige erhaltene Aufzeichnungen scheinen darauf hinzudeuten, dass \textsc{Gauss} zunächst die Klassen für diejenigen Determinanten berechnet hat, die derselben Hundertergruppe und demselben Rest beim Teiler 15 angehören. Die Determinanten dieser Abteilungen wurden dann nach der Anzahl der Geschlechter und Klassen geordnet, zuletzt alle derselben Hundertergruppe angehörigen auf die hier wiedergegebene Weise.
In die Tafeln der einzelnen Hundertergruppen sind manche spätere Berichtigungen eingefügt, nicht aber in die Zusammenstellungen in Tausendergruppen. Datierungen enthalten nur die beiden Tafeln mit den Determinanten der Formen $-(15n+7)$ und $-(15n+13)$, und zwar entsprechend ‘\textit{Expl. In. Febr. 1801}’ und ‘\textit{Expl. 27 Febr. 1807.}’
\par In diesen Tafeln habe ich unter anderem die folgenden Fehler bemerkt. Zur leichteren Kontrolle füge ich die Periodenzahlen der Fundamentalklassen hinzu, etwa $4\pdot\ 4\pdot\ 2$ bei der Determinante
